% =====================================================================
\chapter{Autoevaluación}
\label{ap:autoevaluacion}
% =====================================================================

\section{Puntuación por parte y por módulo}

El examen se corrige con la clave del \cref{tab:clave} y con los criterios entre corchetes de cada solución. El \cref{tab:puntos-parte} recoge la puntuación por parte y el \cref{tab:puntos-modulo}, la puntuación por módulo, que sirve para localizar las debilidades.

\begin{table}[h]
\centering\small
\caption{Puntuación por parte.}\label{tab:puntos-parte}
\begin{tabular}{@{}lcc@{}}
\toprule
Parte & Puntos & Nota obtenida \\
\midrule
A. Cuestiones (20 de 2 puntos) & 40 & \\
B. Problemas (6 de 7 puntos y 3 de 6 puntos) & 60 & \\
\midrule
\textbf{Examen} & \textbf{100} & \\
\midrule
C. Proyecto de fase (se califica aparte) & 100 & \\
\bottomrule
\end{tabular}
\end{table}

\begin{table}[h]
\centering\small
\caption{Puntuación por módulo.}\label{tab:puntos-modulo}
\begin{tabularx}{\textwidth}{@{}>{\raggedright\arraybackslash}p{4.6cm}>{\raggedright\arraybackslash}X>{\raggedright\arraybackslash}Xccc@{}}
\toprule
Módulo & Cuestiones & Problema & Puntos & Mínimo & Obtenidos \\
\midrule
0.1. Lógica y demostraciones & 1, 2 (4~p) & B.1 (7~p) & 11 & 8 & \\
0.2. Conjuntos y funciones & 3, 4 (4~p) & B.2 (7~p) & 11 & 8 & \\
0.3. Números, álgebra y sumatorios & 5, 6, 7 (6~p) & B.3 (6~p) & 12 & 9 & \\
0.4. Trigonometría y complejos & 8, 9 (4~p) & B.4 (7~p) & 11 & 8 & \\
0.5. Cálculo diferencial & 10, 11, 12 (6~p) & B.5 (7~p) & 13 & 10 & \\
0.6. Cálculo integral y series & 13, 14 (4~p) & B.6 (6~p) & 10 & 7 & \\
0.7. Vectores y varias variables & 15, 16 (4~p) & B.7 (7~p) & 11 & 8 & \\
0.8. Combinatoria y probabilidad & 17, 18 (4~p) & B.8 (6~p) & 10 & 7 & \\
0.9. Estadística & 19, 20 (4~p) & B.9 (7~p) & 11 & 8 & \\
\midrule
\textbf{Total} & 20 cuestiones (40~p) & 9 problemas (60~p) & \textbf{100} & 70 & \\
\bottomrule
\end{tabularx}
\end{table}

La columna «Mínimo» es el 70\,\% de los puntos del módulo, redondeado hacia arriba. Una puntuación inferior indica que conviene repasar ese módulo aunque el examen esté aprobado.

\section{Umbral de aprobado}

\begin{itemize}[label=\textendash]
  \item \textbf{Examen de hito}: 70 puntos de 100 (70\,\%) o más.
  \item \textbf{Proyecto de fase}: 70 puntos de 100 o más en la rúbrica del \cref{tab:rubrica}.
  \item \textbf{Puerta de salida de la fase~0}: además del examen, cada uno de los módulos 0.1 a 0.9 exige un test de módulo con al menos 9 respuestas correctas de 12, y la autoevaluación del laboratorio interactivo de la fase debe alcanzar al menos el 70\,\%.
\end{itemize}

\section{Qué hacer si no se alcanza el umbral}

\begin{enumerate}[label=\arabic*.]
  \item Anotar cada fallo en el registro de errores y clasificarlo como error de concepto, de cálculo o despiste.
  \item Repasar las secciones del \cref{tab:repaso} correspondientes a los fallos, empezando por los errores de concepto.
  \item Rehacer sin mirar la solución los ejercicios de nivel $\star\star$ y $\star\star\star$ de esos temas en los apuntes del módulo. En los fallos de demostración, redactar de nuevo la prueba completa y contrastarla con las pautas de redacción del Módulo 0.1, §6.4.
  \item Repetir el test del módulo afectado hasta alcanzar 9 de 12 y volver a presentarse al examen de hito con una versión nueva, no antes de una semana.
\end{enumerate}

\begin{table}[h]
\centering\small
\caption{Secciones de repaso según el fallo.}\label{tab:repaso}
\begin{tabularx}{\textwidth}{@{}>{\raggedright\arraybackslash}p{3.6cm}X@{}}
\toprule
Fallo & Repasar \\
\midrule
Cuestiones 1, 2; problema B.1 & Módulo 0.1, §1.2 (recíproco y contrarrecíproco), §1.4 (equivalencia lógica), §2.4 (negación de cuantificadores), §3.3 (contrarrecíproco), §5.1 y §5.4 (inducción y definiciones recursivas) \\
Cuestiones 3, 4; problema B.2 & Módulo 0.2, §2.3 (imagen y preimagen), §3.1 (inyectiva y sobreyectiva), §4.2 (composición) y §4.4 (cálculo de inversas) \\
Cuestiones 5, 6, 7; problema B.3 & Módulo 0.3, §2.4 (medias), §4.1 y §4.3 (inducción y sumas clásicas), §5.2 y §5.5 (logaritmos, bits y nats), §6.2 y §6.3 (binario y complemento a dos) \\
Cuestiones 8, 9; problema B.4 & Módulo 0.4, §2.2 y §2.3 (suma, diferencia y ángulo doble), §6.2 (forma polar), §6.4 (De Moivre) y §6.5 (raíces $n$-ésimas y de la unidad) \\
Cuestiones 10, 11, 12; problema B.5 & Módulo 0.5, §2.3 y §2.4 (límites), §4.1 y §4.3 (derivada), §5.2 (teorema del valor medio), §5.3 (L'Hôpital) y §6.1 (polinomio de Taylor) \\
Cuestiones 13, 14; problema B.6 & Módulo 0.6, §3.1 y §3.2 (cambio de variable y partes), §4.1 (integrales impropias), §5.3 (series $p$) y §6.1 a §6.4 (series de potencias y de Taylor) \\
Cuestiones 15, 16; problema B.7 & Módulo 0.7, §4.1, §4.2 y §4.4 (parciales, plano tangente y gradiente), §5.2 (hessiano), §5.3 y §5.4 (Weierstrass y Lagrange) y §6.3 (polares) \\
Cuestiones 17, 18; problema B.8 & Módulo 0.8, §1.4, §2.1 y §2.2 (variaciones, combinaciones y multinomiales), §2.4 (doble conteo), §4.3 (Bayes) y §5.2 y §5.4 (esperanza y binomial) \\
Cuestiones 19, 20; problema B.9 & Módulo 0.9, §1.4 (desviación típica), §2.5 (varianza muestral), §4.2 (máxima verosimilitud), §4.3 y §4.4 (intervalos), §5.1 a §5.3 (contrastes) y §5.6 (comparaciones múltiples) \\
\bottomrule
\end{tabularx}
\end{table}

\paragraph{Siguiente paso.} Superado el hito, se continúa con la fase~I, módulo~1: modelos de lenguaje probabilísticos.
